\subsection{Summary of Document Layout Analysis}
\subsubsection{Geometric Structure Detection and Analysis}

This line of work treats document layout as a purely visual or geometric structure, with the core objective of locating and classifying elements such as text blocks, images, and tables on a page. These methods output structured layout segmentation results that serve as fundamental inputs for subsequent document understanding tasks. Early work, such as The Document Spectrum for Page Layout Analysis \cite{ogorman1993documentspectrum}, characterized the overall page structure from a signal processing perspective. Document layout analysis for semantic information extraction \cite{ferrara2017documentlayoutsemantic} and Visual Detection with Context for Document Layout Analysis \cite{soto2019visual} gradually combined layout analysis with information extraction requirements and contextual modeling; the latter introduced object detection frameworks like Faster R-CNN and explicitly utilized contextual information to improve detection robustness in complex layouts. With the development of deep learning, a large body of research began adopting CNN- or Transformer-based detection models for layout analysis. For instance, Document Layout Analysis with an Enhanced Object Detector \cite{minouei2021enhancedobjectdetector} used convolutional neural networks to estimate document structure boundaries, whereas Transformer-Based Approach for Document Layout Understanding \cite{yang2022transformerlayout} introduced global self-attention for document layout understanding. Document Layout Analysis Via Positional Encoding \cite{zhou2022positionalencodinglayout} employed positional encoding to alleviate the insufficient modeling of long-range spatial relationships by traditional detectors in complex layouts. VSR: A Unified Framework for Document Layout Analysis Combining Vision, Semantics and Relations \cite{zhang2021vsr} enhanced interaction modeling between layout elements by fusing vision, semantics, and relation modules. a Graphical Approach to Document Layout Analysis \cite{wang2023graphicallayout} utilized graph structures to establish global relationships between PDF elements and improve efficiency. Advanced Layout Analysis Models for Docling \cite{livathinos2025docling} evaluated the practicality of various object detection architectures, such as RT-DETR and DFINE, in industrial-grade systems. Addressing issues of generalization and adaptation to real-world scenarios, In-domain versus Out-of-domain Transfer Learning for Document Layout Analysis \cite{li2024indomain} systematically analyzed cross-domain performance degradation, while LD-DOC proposed a lightweight domain-adaptive layout analysis method to enhance cross-domain generalization \cite{gao2024lddoc}. Regarding detection framework design, recent methods have continuously improved model unification and efficiency. a Hybrid Approach for Document Layout Analysis in Document Images \cite{hybrid_dla_2024} enhanced the multi-scale and multi-category recognition capabilities of Transformer-based detectors through a hybrid matching mechanism. DocLayout-YOLO \cite{doclayout_yolo_2024} balanced speed and accuracy by combining high-quality synthetic data with global-to-local perception modules. PP-DocLayout \cite{pp_doclayout_2025} proposed a unified detection model to support large-scale data construction, while DLAFormer \cite{dlaformer_2024} and Doc-DINO \cite{doc_dino_2024} further explored end-to-end Transformer architectures to unify the modeling of detection, classification, and logical structural relationships.

\subsubsection{Layout-Aware Representation Learning}

In this stage, layout information is explicitly integrated into the model’s representation learning process and jointly modeled with text content, enabling the model to capture implicit spatial structures and semantic relationships at the representational level. Since LayoutLM \cite{xu2019layoutlm} first introduced the joint modeling of text content and 2D positional information into the Transformer pre-training framework, document layout understanding transitioned from pure visual parsing to a representation learning paradigm. This work allowed the model to jointly learn text and layout during pre-training by modeling token coordinate embeddings alongside text embeddings. However, relying solely on coordinate embeddings and text masking tasks made it difficult to fully characterize complex spatial relationships. Position Masking \cite{saha2021position} proposed masking both text and positional embeddings during pre-training, forcing the model to explicitly predict spatial positions and thereby strengthening its ability to model positional semantics.

In the multimodal direction, LAMPRET \cite{wu2021lampret} further introduced visual features into layout-aware pre-training, aligning text, layout, and visual information through cross-modal contrastive learning. Addressing the issue that traditional methods implicitly rely on language-specific features, LiLT \cite{wang2022lilt} proposed a language-independent Layout Transformer that modeled spatial relationships between tokens through a layout-only self-attention mechanism, enhancing robustness in cross-lingual and low-resource scenarios. XYLayoutLM \cite{gu2022xylayoutlm} focused on OCR ordering noise and the insufficient modeling of reading orders, introducing the Augmented XY Cut algorithm to generate more reasonable token sequences combined with dilated conditional position encodings to improve stability in complex layouts. Regarding the expressive power of spatial relationships, ERNIE-Layout \cite{peng2022ernie_layout} explicitly modeled richer spatial and structural relationships on a page through layout-knowledge-enhanced pre-training tasks. Beyond Layout Embedding \cite{zhu2023beyond_layout} approached the problem from the attention mechanism level, introducing Gaussian-bias-based layout attention, allowing the model to dynamically adjust attention distributions based on the relative distance and orientation between entities.

Furthermore, LayoutMask \cite{tu2023layoutmask} strengthened the semantic interaction between text and layout during pre-training by designing differentiated masking strategies for different layout units. Beyond pre-training objectives and attention mechanisms, the layout representation format itself has continued to evolve. Multi-scale Cell-based Layout Representation \cite{shi2023multiscale_cell} proposed a cell-level layout representation based on row and column indices, fusing word-level and cell-level spatial information through a multi-scale design to make layout expression closer to the document’s hierarchical logical structure. In recent years, research has begun to explicitly model layout structures in relational forms. Modeling Layout Reading Order as Ordering Relations \cite{zhang2024reading_order_relations} elevated reading order from a linear sequence to directed ordering relations, injecting these relations as enhancement signals into the layout-aware model’s attention. Based on this, a Graph-Augmented Multi-Stage Transformer Model \cite{arshad2025graph_augmented} further introduced graph neural networks and a multi-stage Transformer architecture to progressively refine fine-grained structural representations by explicitly modeling the spatial and semantic relationships between layout components.

\subsubsection{Layout-Aware Reasoning and Generation}

As the capabilities of Large Language Models (LLMs) have improved, researchers have gradually realized that relying solely on text content is insufficient for complex document understanding; the layout structure of a document provides critical reasoning cues. Although early research modeled layout as an input feature, the robustness of such approaches in real-world scenarios remained limited. Unveiling the Deficiencies of Pre-trained Text-and-Layout Models in Real-world Visually-rich Document Information Extraction \cite{zhang2024unveiling} systematically analyzed the performance degradation of models like LayoutLM and ERNIE-Layout on real documents, pointing out that existing methods struggled to fully utilize complex layout structures. This has pushed research from treating "layout as an input feature" toward "layout participating in the reasoning process." Based on this insight, DocLLM \cite{wang2024docllm} allowed layout information to directly influence the model's attention calculation by explicitly modeling the spatial positions of text blocks, thereby enhancing generation capabilities for structurally complex documents. mPLUG-DocOwl 1.5 \cite{hu2024mplugdocowl15} further unified the modeling of document structure from a visual perspective, enabling the model to understand layout relationships without relying on OCR. Meanwhile, LayoutLLM \cite{luo2024layoutllm} enabled the model to actively consider layout structure during generation through layout instruction tuning. LAPDoc \cite{lamott2024lapdoc} explicitly incorporated document structure into prompts through layout-aware prompting, reducing reasoning errors caused by ignoring layout. Furthermore, ReLayout \cite{relayout2024} focused on the lack of precise annotations in real-world applications, enabling the model to learn effective structural information relying only on OCR through layout-enhanced pre-training. Collectively, these works indicate that layout has transitioned from a simple input feature to an important condition affecting the model's reasoning process.

As task complexity increased, researchers further discovered that in cross-page and multi-hop reasoning scenarios, using layout only to assist reasoning is still insufficient because key information is often distributed across different locations. Consequently, layout began to be used to organize the retrieval process, helping models locate relevant evidence more accurately. In this direction, LAD-RAG \cite{ladrag2025} enabled models to perform cross-page retrieval using layout structures by constructing document layout graphs. SuperRAG \cite{superrag2025} similarly modeled layout relations as graph structures and combined them with graph neural networks to improve multi-hop reasoning capabilities. Structured Multimodal RAG \cite{structuredmmrag2025} improved the relevance of retrieval results and reduced erroneous generation by utilizing spatial relationships between document blocks. Meanwhile, research has shown that the method of partitioning retrieval units significantly impacts performance. Document Segmentation Matters for RAG \cite{docseg2025} proposed a pseudo-instruction-based chunking method to make retrieval blocks more aligned with semantic structures. Enhancing RAG System Performance Through Semantic Layout Chunking \cite{semanticlayoutchunking2025} combined semantic and layout information for chunking to maintain document structural consistency. In more structured document scenarios, TabRAG \cite{tabrag2025} utilized the row-column structure of tables to improve retrieval effectiveness, while From Chart to QA Pairs \cite{chart2qapairs2025} constructed a QA dataset for chart-based documents, driving the development of layout-aware reasoning research from a data perspective. These methods demonstrate that layout can serve not only as an auxiliary for reasoning but also as a key structure for organizing the retrieval process.

Building upon these retrieval-oriented approaches, research has further explored treating layout itself as a generation target, enabling models to actively generate document layouts that conform to structural constraints. LayoutKAG \cite{layoutkag2024} improved the stability of layout generation by combining layout knowledge with retrieval examples. LayoutCoT \cite{shi2025layoutcot} decomposed layout generation into a step-by-step reasoning process, resulting in more structurally consistent layouts. LayoutRAG \cite{layoutrag2025} addressed information deficiency by retrieving relevant layout templates to assist generation. CAL-RAG \cite{calrag2025} enabled the coordination of content and layout generation through multi-agent collaboration. OmniDocLayout \cite{omnidoclayout2025} allowed the model to generate an overall structure before localized details through a step-by-step refinement approach. These works mark the transformation of layout from a passive input to an object of active model generation.

At an even higher level of integration, the latest research has begun to treat layout as a unified structure connecting the retrieval, reasoning, and generation modules, making it a core component of a complete system. CMRAG \cite{cmrag2025} made the retrieval and generation processes more consistent by jointly modeling visual, textual, and layout information. SymbioticRAG \cite{symbioticrag2025} introduced human feedback, allowing layout information to guide retrieval and generation at critical steps. HySemRAG \cite{hysemrag2025} improved the stability and accuracy of complex document reasoning by combining layout structure with semantic information.

Taken together, these developments reveal a clear trajectory: the role of layout has evolved from an input feature to a reasoning condition, a retrieval structure, a generation target, and finally, a core structure connecting all modules, playing an increasingly critical role in complex document understanding and generation.

\subsubsection{Layout Analysis Evaluation}

The development of the document layout analysis and understanding field exhibits a progressive evolution from task-driven to data-driven approaches, and from automated labeling to fine-grained human annotation. In early explorations, A Region-based Document VQA \cite{jiang2022regionbased} attempted to integrate information from different document regions into QA models, allowing models to reason by combining layout and text content, which laid the foundation for the application of visual-spatial understanding in document QA. Shortly thereafter, the DocVQA \cite{mathew2020docvqa} dataset constructed document visual question-answering tasks through large-scale question-document pairs, pushing document understanding from simple text recognition toward QA reasoning. In the field of layout analysis, PubLayNet \cite{zhong2019publaynet} achieved layout annotation for over 360,000 pages by automatically aligning PDFs and XML files of academic articles, providing a reliable training foundation for deep learning models. Subsequently, DocBank \cite{li2020docbank} further refined layout element labels, enabling models to understand document structures across richer dimensions. To enhance model generalization across multi-domain documents, DocLayNet \cite{pfitzmann2022doclaynet} provided over 80,000 pages of fine-grained human-annotated data from various fields, allowing layout analysis to cover diverse scenarios such as academic, financial, patent, and legal documents. Meanwhile, the ICDAR competition series has long released various layout analysis tasks, from early complex layout recognition competitions to the latest ICDAR 2023 Competition on Robust Layout Segmentation in Corporate Documents \cite{auer2023icdar23}, driving the community to continuously improve algorithmic robustness and practicality under unified evaluation mechanisms. Overall, these representative works and related competition datasets constitute the core pillars of document layout analysis research, illustrating an evolutionary trajectory from visual question answering and automated labeling to diversified human annotation, and finally to comprehensive evaluation and robustness analysis.
